\chapter{Floating Point Arithmetic}\label{capitolo:2}
This chapter will be devoted to floating point arithmetic. It will start with a refresher of the well known basic principles, with the aim of presenting how arbitrary precision computation works in Julia.  

\section{Floating point number system}
\textcolor{red}{La solita spiegazione. Si può guardare~\cite[§~2.1]{higham:accuracy_book}}

\textcolor{red}{Cose da dire: significand/mantissa. Cos'è la precisione di macchina e cosa l'unit roundoff (è importante?), introdurre la notazione $u_{53}$.}

\section{Multiple precision arithmetic}
\textcolor{red}{Qualcosa si trova in~\cite[§~27.9]{higham:matfun_book}.}
\smallskip 

\textcolor{red}{Uhm che incipit faccio?}
\smallskip 

\textcolor{red}{Cose da dire: precision e accuracy non sono la stessa cosa.}

\subsection{Multiple precision arithmetic in \textsc{matlab}}
In \matlab{}, arbitrary precision arithmetic is commonly performed via the \texttt{vpa} command. The commercial Multiprecision Computing Toolbox~\cite{matlab_mpct} provides mathematical functions that natively support the toolbox's multiprecision floating-point numeric type.

\textcolor{red}{Sarebbe anche interessante fare un breve confronto tra vpa e Advanpix (magari anche solo riassumendo quanto sta nella documentazione (per vpa), e linkando alla pagina del manuale di Advanpix). Però, se la faccio, la farei breve perché alla fine MATLAB manco lo uso. Se questo confronto/approfondimento fa emergere degli aspetti generali dell'aritmetica floating point, allora può aver senso.}

\subsection{Multiple precision arithmetic in Julia}
The Julia programming language~\cite{article:Julia} natively supports multiprecision floating-point numbers by means of the built-in data type \texttt{BigFloat}.

\texttt{BigFloat}s support standard arithmetic operations as well as many mathematical functions, each of which is implemented via a corresponding call to the \textsc{mpfr}~\cite{article:MPFR} library. In addition, \texttt{BigFloat}s provide a set of dedicated methods, some of which are described in the following sections.

\subsubsection*{Basic functioning of \texttt{BigFloat}s}
\textcolor{red}{Forse un cappello introduttivo su come sono fatti i BigFloats grosso modo (uno struct etc), però non sono sicurissimo neanche io. Scavare in quel codice non è facile.}
\smallskip 

\textcolor{red}{Avevo fatto qualche benchmark, in cui è emerso come il costo delle operazioni tra \texttt{BigFloat}s dipende dalla precisione (che ci si aspetta d'altronde):  maggiore è la precisione, maggiore il costo in termini di tempo e di memoria utilizzata (ovviamente). Inoltre il numero di allocazioni cresce, ma contrariamente alle due statistiche precedenti, la dipendenza non è lineare, bensì si assiste a un ``salto''. Magari è interessante metterlo qui. Però ho testato con le matmuls, quindi lo metto dopo? ripeto il test con i prodotti ``normali''?}

\subsubsection*{Common pitfalls in multiprecision arithmetic}
The \texttt{precision} and \texttt{setprecision} commands are well explained in the language documentation~\cite{link:julia_documentation}.
Examples~\ref{es:setprecision_global} and~\ref{es:setprecision_next_ops} highlight two important aspects.

\begin{esempio}\label{es:setprecision_global}
	%The \texttt{setprecision} command changes the default precision globally.
	The $\mathtt{setprecision}(\mathtt{T}, p)$ call sets the precision of type \texttt{T} to $p$ \emph{globally}. The \texttt{setprecision} command mutates the global state, its effects are visible even outside the block in which the call happens.
\begin{lstlisting}[language=Julia]
  # precision(BigFloat) = 256 (default)
  a = 74	
  let 
    a = 42
    setprecision(BigFloat, 512)
  end 
  # a = 74, precision(BigFloat) = 512
\end{lstlisting}
	$\mathtt{setprecision}(p)\ \mathtt{do}\ \dots\ \mathtt{end}$ changes the precision to $p$ for the duration of the evaluation of the \texttt{do} block.
\end{esempio}

\begin{esempio}\label{es:setprecision_next_ops}
	The \texttt{setprecision} command doesn't retroactively change the precision of already performed computations. Instead, the effects of this command only affect the subsequent \emph{computations} throughout the program.
\begin{lstlisting}[language=Julia]
  # In Julia, π is a mathematical constant
  pi = BigFloat(π) # pi = 3.1415...6198 
        	   # (78 correct decimal digits)
  setprecision(BigFloat, 53)
  pi 		# still displays 78 decimal digits 
  pi * 1.0	# 3.1415926535897931
  precision(pi)	# 256
\end{lstlisting}
\end{esempio}
\smallskip 

\noindent Example~\ref{es:big_BigFloat_constructor} llustrate some subtleties in the usage of the \texttt{BigFloat()} and \texttt{big()} constructors.
\begin{esempio}\label{es:big_BigFloat_constructor}
	Let $\mathtt{x}$ be a variable or numerical expression of type \texttt{BigFloat}. The $\mathtt{BigFloat}(\mathtt{x})$ and $\mathtt{big}(\mathtt{x})$ constructors work slightly differently.
	The former returns a value with the precision set to the current working precision; instead $\mathtt{big}(\mathtt{x})$ returns a maximum precision representation of $\mathtt{x}$.
\begin{lstlisting}[language=Julia]
  pi = BigFloat(π)
  setprecision(53) do 
    big(pi)	 # 3.1415...6198 (78 decimal digits)
    BigFloat(pi) # 3.1415926535897931
  end
\end{lstlisting}
\end{esempio}
\smallskip

\noindent The uncareful usage of the \texttt{big} and \texttt{BigFloat} commands can easily lead to some outcomes that may be unintuitive at first. The main one -- which underlies most of the other pitfalls when converting between double and arbitrary precision -- is illustrated in example~\ref{es:float_promotion}. 
\begin{esempio}\label{es:float_promotion}
	If $\mathtt{x}$ is not of type \texttt{BigFloat}, then the \texttt{big(x)} and \texttt{BigFloat(x)} constructor try to convert it. The conversion of a double precision number is conceptually simple: the significand of $\mathtt{x}$ is preserved, and the extra lower-order significand bits in the high precision representation are set to zero.  
\begin{lstlisting}[language=Julia]
  BigFloat(0.1) # 0.10000000000000000555111512312...
  big(0.1)	# same as BigFloat(0.1)
  big(0.0625)	# 1/16 is exactly representable 
  		# in binary
\end{lstlisting}
\end{esempio}
\smallskip

\noindent Be aware that evaluation happens before promotion, as shown in example~\ref{es:eval_expr}.
\begin{esempio}\label{es:eval_expr}
	When calling \texttt{BigFloat()} or \texttt{big()} on an expression, the expression is evaluated first. 
\begin{lstlisting}[language=Julia]
  big(1/10) 	# only 17 decimal digits are correct 
  1 / big(10)	# all 78 decimal digits are correct
\end{lstlisting}
	In the first case, since the operands are \texttt{Int64}, $1/10$ was evaluated in double precision, as this is Julia's normal behaviour, the result is then promoted. In the second case, $1$ was promoted to a \texttt{BigInt} and a division between arbitrary precision quantities was performed.
\end{esempio}
%\smallskip

\noindent Examples~\ref{es:float_promotion} and~\ref{es:eval_expr} show that neither \texttt{big(0.1)} nor $\mathtt{big(1/10)}$ is either the high precision evaluation of the expression $\mathtt{1/10}$ or the closest \texttt{BigFloat} to the real number $10^{-1}$. 
The latter can be obtained, in Julia, with \texttt{big(``0.1'')} (or with \texttt{BigFloat(``0.1'')}); but \texttt{big(``1/10'')} and \texttt{BigFloat(``1/10'')} fail.
\smallskip 

\noindent Example~\ref{es:convert} shows an aspect of the \texttt{convert} command's behaviour that can have a concrete impact in multiprecision computations.
\begin{esempio}\label{es:convert}
	The call $\mathtt{convert}(\mathtt{T},x)$ returns a copy unless $x$ is already of type $\mathtt{T}$, in which case an \emph{alias} of $x$ will be returned.
\begin{lstlisting}[language=Julia]
  A = rand(BigFloat, 2,2)
  B = setprecision(53) do 
    convert(Matrix{BigFloat}, A)
  end
  precision(B[1,1])	# 256
  A === B 		# true 
  pointer(A) == pointer(B) # true
\end{lstlisting}
\end{esempio}
Much more on conversion and promotion can be found in the Julia documentation~\cite{link:julia_documentation}.

\subsubsection*{Multiprecision linear algebra in Julia}
Let $A,B\in\R{n}{n}$ such that $A, B$ have elements of type \texttt{BigFloat}. The operations $A+B$, $A\cdot B$, $A^{-1}$ and $A^n$ are natively supported for integer $n$. 

\begin{oss*}
	We make a few quick observations on the aforementioned operations:
	\begin{itemize}
		\item $A+B$ broadcasts the $+$ operation onto the single entries of the matrices 
		\item Since \textsc{BLAS}~\cite{article:BLAS} for matrix multiplications are available only for single and double precision, the computation of $C=A\cdot B$ is implemented with the naive three-way for loop, using the $i,j,k$ ordering. Julia tries to exploit Instruction Level Parallelism by using the \texttt{@simd} macro on the innermost loop. $C_{i,j}$ is updated by using fused multiply-add operations, implemented by \textsc{mpfr}
		\item $\mathtt{inv}(A)$ is doable because the $\mathtt{lu}$ command works with $\mathtt{BigFloat}$ matrices out of the box. As before, no \textsc{lapack} routine can be used; the factorization is computed using the gaussian elimination algorithm with pivoting. The lower level arithmetic operations are performed via \textsc{mpfr} calls. 
		\item $A^n,\ n\in\mathbb{Z}$ calls the $\mathtt{power\_by\_squaring}$ function, which essentially is an implementation of the \emph{binary powering} algorithm~\cite[Alg.~4.1]{higham:matfun_book}. For $n<0$, $\mathtt{inv}(A)^{\abs{n}}$ is computed. 
	\end{itemize}
\end{oss*}
If the exponent $n$ is not an integer, $A^n$ is not doable natively. This is because $A^n$ is computed as $Q\,T^n\,T^r\,Q^{h}$, where $A=QTQ^H$ is a Schur decomposition of $A$, $T^r$ with $r=n-\lfloor n\rfloor$ is computed using the algorithm described in~\cite{higham:hili13}.

The Schur decomposition of matrices with \texttt{BigFloat} elements is available through the \texttt{GenericSchur.jl}~\cite{sw:GenericSchur} library, but is not natively implemented in the language.
The library also allows the computation of the spectrum of $A$, or in general its eigendecomposition (namely, it gives the ability to run the Julia \texttt{eigvals} and \texttt{eigen} commands). 

